The analysis of competitive interactions between agents often leads to challenging mathematical optimization problems \cite{daskalakischallenge}. Game theory traditionally assumes that agents are motivated solely by the goal of minimizing convex cost functions \cite{Brooks2023}. This assumption is primarily a mathematical convenience to ensure the tractability of optimization problems and the existence of well-known solution concepts. When this assumption is violated, computing various stable states of interactions becomes difficult, and even determining their existence can be computationally hard \cite{daskalakisstonr}. Similarly, when agents can choose from an infinite number of distinct actions, finding actions that result in favorable outcomes reduces to global optimization over nonconvex functions.

Nevertheless, there are combinations of reasonable solution concepts and classes of games that are both tractable and widely applicable. Allowing agents to play randomized strategies guarantees the existence of stable points but still poses the challenge of global optimization. In this chapter, we introduce two prominent approaches to global optimization: Lasserre's hierarchies and the spatial Branch and Bound method. In later chapters, these methods will enable us to solve games with infinite action spaces and nonconvex cost functions. Furthermore, we will demonstrate their utility in the context of robotics.